\chapter{The Gaussian Chain by Belief Propagation}\label{ch:gaussian-bp}

Chapter~\ref{ch:gaussian} obtained the joint score by inverting a matrix. This
chapter obtains the same score by passing messages along the chain, which costs
$O(L)$ updates instead of an $O(L^3)$ inversion and, unlike the matrix route,
survives the loss of Gaussianity in Chapter~\ref{ch:laplace}. Three things
follow: the recursion is exactly the Kalman filter and the
Rauch--Tung--Striebel smoother, update for update; the posterior has a closed
form in the bulk; and the influence of distant evidence decays geometrically,
which is what makes a local approximation of the score possible at all.

\section{Belief propagation on the chain}\label{sec:bp-chain}\label{sec:bp-problem}

The matrix route inverts an $L \times L$ matrix, a global computation
over the whole sequence. The score--posterior identity says the same object is a posterior
expectation, and the posterior has a structure that supports local
computation. This section develops that computation once, in full.

\subsection{The factor graph of the posterior}

For any Markov-chain prior and per-frame OU evidence, Bayes' rule gives
\begin{equation}
  p(a \,|\, x) \;\propto\;
  \psi_{\mathrm{pr}}(a_0)
  \prod_{k=0}^{L-2} \psi^{(k)}_{\mathrm{tr}}(a_k, a_{k+1})
  \prod_{k=0}^{L-1} \psi^{(k)}_{\mathrm{ob}}(a_k),
  \qquad
  \psi^{(k)}_{\mathrm{ob}}(a_k) = \Nd\big(x_k;\, e^{-t}a_k,\, \Deltat\big),
  \label{eq:bp-posterior}
\end{equation}
a hidden Markov model \citep{rabiner1989tutorial}: a latent chain observed
through independent per-frame noise. By \eqref{eq:tweedie-joint}, the
$k$-th score component requires exactly the smoothing marginal
$\E[a_k | x]$ of \eqref{eq:bp-posterior}.

The factor graph of this posterior has $L$ variable nodes, $2L$ factor nodes
(one prior, $L{-}1$ transitions, $L$ observations; the data $x_k$ are constants
clamped inside their factors) and $3L - 1$ edges. One fewer edge than nodes, on
a connected graph, is exactly the tree condition of
Section~\ref{sec:factor-graphs}.

\begin{proposition}\label{prop:chain-tree}
For any Markov-chain prior and any per-frame observation model, the
posterior $p(a \,|\, x)$ factorises over a tree-shaped factor graph, and
all $L$ smoothing marginals are computed exactly by one forward and one
backward sweep of sum--product messages.
\end{proposition}

Proposition~\ref{prop:chain-tree} counts sweeps, not work: it says
nothing about how the messages are represented, and for continuous
variables each message is a function
(Section~\ref{sec:sum-product}). Keeping tree-exactness of the
algorithm (a graph property) separate from representability of the
messages (a property of the model's function class) is the conceptual
hinge of this thesis; the Gaussian case closes the gap, and
Chapter~\ref{ch:laplace} measures what happens when it no longer does.

\subsection{The sum--product recursion, with the bookkeeping explicit}

Specialising the sum--product rules \eqref{eq:bp-rules} to this tree,
the schedule collapses into one forward and one backward sweep
(Theorem~\ref{thm:tree-exact}), and the bookkeeping deserves care.

\begin{proposition}[Chain sweeps, Convention A]\label{prop:bp-convA}
Define messages
\begin{align}
  \mu_{\to 0}(a_0) &:= \psi_{\mathrm{pr}}(a_0),
  &
  \mu_{\to k+1}(a_{k+1}) &:= \int
    \psi^{(k)}_{\mathrm{tr}}(a_k, a_{k+1})\,
    \psi^{(k)}_{\mathrm{ob}}(a_k)\,
    \mu_{\to k}(a_k)\; \dd a_k ,
  \label{eq:bp-fwd}\\
  \mu_{\leftarrow L-1}(a_{L-1}) &:= 1,
  &
  \mu_{\leftarrow k-1}(a_{k-1}) &:= \int
    \psi^{(k-1)}_{\mathrm{tr}}(a_{k-1}, a_k)\,
    \psi^{(k)}_{\mathrm{ob}}(a_k)\,
    \mu_{\leftarrow k}(a_k)\; \dd a_k .
  \label{eq:bp-bwd}
\end{align}
Then for every $k$,
\begin{equation}
  p(a_k \,|\, x) \;\propto\;
  \mu_{\to k}(a_k)\;
  \psi^{(k)}_{\mathrm{ob}}(a_k)\;
  \mu_{\leftarrow k}(a_k) :
  \label{eq:bp-combine}
\end{equation}
the forward message into $a_k$ carries the evidence
$x_0, \dots, x_{k-1}$ only, the backward message carries
$x_{k+1}, \dots, x_{L-1}$ only, and the local evidence enters once, at
combination time.
\end{proposition}

The proof is direct marginalisation: integrating \eqref{eq:bp-posterior}
over all $a_j$, $j \neq k$, splits at $a_k$ by the Markov property into
independent left and right halves, and peeling the innermost variable of
each half yields the recursions.

\begin{remark}[The double-counting trap]\label{rmk:bp-trap}
An alternative convention absorbs the local evidence into the forward
message, $\mu^{\mathrm{B}}_{\to k} = \mu_{\to k} \cdot
\psi^{(k)}_{\mathrm{ob}}$; the product \eqref{eq:bp-combine} then counts
$x_k$ twice and must be corrected by dividing one factor back out. The cost
of getting this wrong is not subtle: under
Convention A the BP score matches the matrix score to $10^{-14}$, and under the
naive alternative the disagreement is of order one. Every message in this thesis
is therefore stated with its evidence content explicit.
\end{remark}

\section{The Gaussian closure and the Kalman smoother}
\label{sec:bp-gaussian}\label{sec:bp-kalman}

For the Gaussian chain, two facts close the recursion, and they are the only
ones needed: a product of Gaussians in the same variable is Gaussian
\eqref{eq:gauss-product}, and an affine map of a Gaussian, marginalised, is
Gaussian \eqref{eq:gauss-predict}. Together they say that the recursions
\eqref{eq:bp-fwd}--\eqref{eq:bp-bwd} perform exactly one product and one
affine marginalisation per step. The local evidence, read as a function
of $a_k$, is itself Gaussian,
$\psi^{(k)}_{\mathrm{ob}}(a_k) \propto
\Nd(a_k;\, e^{t} x_k,\, \Deltat e^{2t})$: a pseudo-observation at the
deconvolved location $e^t x_k$ with precision $e^{-2t}/\Deltat$, the
pinning strength that appears throughout this chapter.

\begin{theorem}[Gaussian closure of BP on the chain]\label{thm:bp-closure}
For the Gaussian AR(1) chain under the OU channel, every BP message and
every belief is exactly Gaussian, at every node, for every
$(L, \alpha, t)$. Each message is carried by two numbers, each update
is a few arithmetic operations, and the two sweeps compute the exact
posterior marginals, hence the exact joint score. Here, and only here,
the sweep count of Proposition~\ref{prop:chain-tree} translates into an
actual computation of fixed, small cost per frame.
\end{theorem}

\begin{proof}[Derivation]
Induction along each sweep. The initial messages are Gaussian (the
prior; the flat message is the zero-precision limit of the family). Each
step of \eqref{eq:bp-fwd}--\eqref{eq:bp-bwd} is one application of
the product identity \eqref{eq:gauss-product} (absorb the local evidence) followed
by one of Lemma~2 (push through the transition); the combination
\eqref{eq:bp-combine} is Lemma~1 twice. Exactness follows from tree
exactness of sum--product (Theorem~\ref{thm:tree-exact}).
\end{proof}

\begin{remark}[Closure is algebra, not a central limit theorem]
\label{rmk:bp-no-clt}
The closure just proved holds at node degree two exactly as it would at
any degree: it is algebra, not an aggregation effect, as discussed in
\compendium{}. The central-limit argument matters only
for the per-node closures of statistical physics
(recorded in \compendium{}), and on the chain those fail in a way that
can be made exact, with a breakdown time $t_c(\alpha)$ and a critical
coupling $\alpha_c = \sqrt{2} - 1$ (\compendium{}).
\end{remark}

Executing the Gaussian sweeps in mean--variance arithmetic reproduces,
update by update, the classical algorithms of linear estimation.

\begin{toolbox}{The Kalman filter, derived from the BP sweep}
Write the forward message in Convention A as
$\mu_{\to k} = \Nd(a_k;\, m_{k|k-1},\, P_{k|k-1})$, the predicted state
given past observations. One BP step does two things.

\emph{Update (absorb the evidence; Lemma 1).} Multiplying by the
pseudo-observation $\Nd(a_k; e^t x_k, \Deltat e^{2t})$ gives the filtered
Gaussian $\Nd(a_k; m_{k|k}, P_{k|k})$ with
\begin{align*}
  \frac{1}{P_{k|k}} &= \frac{1}{P_{k|k-1}} + \frac{e^{-2t}}{\Deltat},\\
  m_{k|k} &= P_{k|k} \left( \frac{m_{k|k-1}}{P_{k|k-1}}
        + \frac{e^{-t} x_k}{\Deltat} \right)
  = m_{k|k-1} + G_k \big( e^{t} x_k - m_{k|k-1} \big),
\end{align*}
where $G_k = P_{k|k}\, e^{-2t}/\Deltat \in (0,1)$ is the Kalman gain: the
filtered mean is the prediction corrected by a gain-weighted innovation
\citep{kalman1960new}.

\emph{Predict (push through the dynamics; Lemma 2).}
\begin{equation*}
  m_{k+1|k} = \alpha\, m_{k|k},
  \qquad
  P_{k+1|k} = \alpha^2 P_{k|k} + \sigma_\eta^2 .
\end{equation*}
The forward BP sweep is the Kalman filter, message by message.

\emph{Smoothing (combine both directions).} The belief
\eqref{eq:bp-combine} multiplies the filtered forward Gaussian by the
backward message (Lemma~1 once more), producing the posterior marginal
given all observations, $\Nd(a_k; m_{k|L-1}, P_{k|L-1})$; carrying out
the same algebra for the backward recursion yields the fixed-interval
smoother of \citet{rauch1965maximum}. The two-sweep BP schedule on the
chain is the Kalman filter followed by the RTS smoother, exactly as the
generic correspondence recorded in \compendium{} predicts.
\hfill$\square$
\end{toolbox}

Combining with the score--posterior identity answers RQ2: the $k$-th
component of the joint diffusion score is an affine function of the
Kalman-smoothed estimate of frame $k$,
\begin{equation}
  S_k(x, t)
  \;=\; \frac{e^{-t}\, m_{k|L-1} - x_k}{\Deltat} .
  \label{eq:bp-score-kalman}
\end{equation}
For Markov sequences, score computation is a smoothing problem, and the
classical theory of linear smoothing applies directly.

\paragraph{Numerical agreement of the two routes.} The audit suite runs BP
against the matrix route over $80$ configurations
($L \in \{2,3,5,8,16\}$, $\alpha \in \{0.2, 0.5, 0.9, -0.4\}$,
$t \in \{0.05, 0.3, 1, 3\}$, random $x$): maximum disagreement
$1.3 \times 10^{-15}$ on posterior means, $3.9 \times 10^{-15}$ on
variances, $1.2 \times 10^{-14}$ on scores, i.e.\ double-precision
rounding accumulated along the recursion. An independent implementation,
written from scratch against the recursion specification alone,
reproduces the same agreement (Appendix~\ref{app:repro}). One qualitative
detail from a worked $L = 3$ instance: at $x = (1.2, -0.4, 0.8)$ the score
component $S_1$ is positive although $x_1 < 0$, because the neighbours
pull the posterior mean of the middle frame positive. No per-frame
marginal score can produce this; it is inter-frame evidence sharing in
numbers.

\section{The bulk in closed form, and the price of locality}
\label{sec:g-bulk}\label{sec:g-locality}\label{sec:bp-truncation}

For a long chain the interior (bulk) is translation invariant and can be
solved exactly. By Proposition~\ref{prop:g-posterior} the posterior
precision in the bulk is tridiagonal Toeplitz with
\begin{equation}
  \boxed{\;
  J_d \;=\; \frac{e^{-2t}}{\Deltat} + \frac{1+\alpha^2}{\sigma_\eta^2},
  \qquad
  \beta \;=\; -\frac{\alpha}{\sigma_\eta^2},
  \qquad \sigma_\eta^2 = 1 - \alpha^2 .
  \;}
  \label{eq:g-bulk-params}
\end{equation}
The diagonal $J_d$ is the sum of the evidence precision $e^{-2t}/\Deltat$
(divergent as $1/2t$ at small $t$, vanishing at large $t$) and the
constant prior precision; the coupling $\beta$ is pure prior. The
dimensionless ratio $|\beta|/J_d \in (0, \tfrac{1}{2})$ controls
everything below.

\begin{theorem}[Exact bulk variance and geometric covariance]
\label{thm:g-bulk-cov}
In the bulk of a long chain, the posterior variance and covariance are
\begin{equation}
  \boxed{\;
  V \;=\; \frac{1}{\sqrt{J_d^2 - 4\beta^2}},
  \qquad
  \big( J^{-1} \big)_{i,\, i+d} \;=\; q^{\,d}\; V,
  \qquad
  q \;=\; \frac{J_d - \sqrt{J_d^2 - 4\beta^2}}{2\,|\beta|} \in (0, 1).
  \;}
  \label{eq:g-bulk-cov}
\end{equation}
\end{theorem}

\begin{toolbox}{Transfer-matrix inversion of a tridiagonal Toeplitz matrix}
Seek the inverse row in geometric form,
$(J^{-1})_{i, i+d} = C q^{\,d}$ for $d \geq 0$, symmetric in $\pm d$. The
defining relation $\sum_j J_{ij} (J^{-1})_{j, i+d} = \delta_{0d}$ reads,
for $d \geq 1$,
\begin{equation*}
  \beta\, C q^{\,d-1} + J_d\, C q^{\,d} + \beta\, C q^{\,d+1} = 0
  \quad\Longleftrightarrow\quad
  \beta q^2 + J_d\, q + \beta = 0,
\end{equation*}
a quadratic whose root inside the unit disc is
$q = (J_d - \sqrt{J_d^2 - 4\beta^2})/(2|\beta|)$ (for $\alpha > 0$ all
inverse entries are positive; for $\alpha < 0$ they alternate in sign with
the same modulus). The $d = 0$ relation fixes the prefactor:
$J_d C + 2\beta C q = 1$, and since
$2\beta q = -(J_d - \sqrt{J_d^2 - 4\beta^2})$,
$C = 1/\sqrt{J_d^2 - 4\beta^2} = V$. Existence of the real square root
requires $J_d > 2|\beta|$, which always holds here:
\begin{equation*}
  J_d - 2|\beta|
  = \frac{e^{-2t}}{\Deltat} + \frac{1 + \alpha^2 - 2|\alpha|}{\sigma_\eta^2}
  = \frac{e^{-2t}}{\Deltat} + \frac{(1 - |\alpha|)^2}{\sigma_\eta^2}
  \;>\; 0
\end{equation*}
for every $|\alpha| < 1$ and $t > 0$. This positivity margin is also what
guarantees the belief-propagation fixed point exists for every coupling
and diffusion time (\compendium{}). \hfill$\square$
\end{toolbox}

The number $q(\alpha, t)$ measures how far evidence travels along the
chain: the posterior correlation length is $\xi = 1/\log(1/q)$. Its limits
are
\begin{equation}
  \begin{aligned}
  t \to 0:&\quad q \approx \frac{|\beta|}{J_d} \to 0
  &&(\text{each frame pinned by its own observation});\\
  t \to \infty:&\quad q \to |\alpha|
  &&(\text{posterior} \to \text{prior}),
  \end{aligned}
  \label{eq:g-q-limits}
\end{equation}
the second because
$q \to \big((1+\alpha^2) - (1-\alpha^2)\big)/(2|\alpha|) = |\alpha|$ when
the evidence term vanishes. As diffusion time runs, the posterior
correlation length climbs monotonically from $0$ to the prior's own
$1/\log(1/|\alpha|)$.

\subsection{The locality law}

RQ4 asks how rapidly the influence of remote observations decays. Define
the local estimator of radius $r$: the posterior mean of $a_k$ given only
the observations within distance $r$,
$m_k^{(r)} = \E[a_k \,|\, x_{k-r}, \dots, x_{k+r}]$. Two facts make the
comparison clean. Any contiguous window of a stationary AR(1) chain is
itself a stationary AR(1) chain, so $m^{(r)}_k$ is exact on its window;
the only approximation is the truncation of range. And both estimators are
linear in $x$, so the error $m_k - m_k^{(r)} = e^\T x$ has a closed-form
RMS over the data law, $\mathrm{RMS}_r = \sqrt{e^\T \Sigt e}$, a
deterministic number with no sampling involved.

\begin{proposition}[Geometric locality: influence coefficients]
\label{thm:g-locality}
In the bulk, the influence of the observation at distance $d$ on the exact
posterior mean is $(J^{-1})_{k,k\pm d} \propto q(\alpha,t)^{\,d}$, exactly,
by Theorem~\ref{thm:g-bulk-cov}. Truncating the evidence at radius $r$
therefore discards a leading term of order $q^{\,r+1}$.
\end{proposition}

What Proposition~\ref{thm:g-locality} controls is the size of the
\emph{discarded} evidence, and that part is exact. It does not by itself
prove that the radius-$r$ estimator's error inherits the rate: $m^{(r)}_k$
is the posterior mean of a \emph{re-normalised} model on the window, so
truncation changes the boundary conditions as well as removing the leading
coefficient, and bounding the difference between those two effects is not
done here.

The rate is instead confirmed numerically, and closely. Fitted slopes of
$\log\mathrm{RMS}_r$ over $r = 2, \dots, 13$ at the centre of an $L = 121$
chain match $\log q$ within $0.2\%$ at $t \in \{0.1, 0.5, 2.0\}$
($\alpha = 0.8$); see Figure~\ref{fig:g-locality}, where the dashed lines
carry the predicted slope with no fitting. The claim this thesis makes is
therefore: an exact geometric law for the influence coefficients, and a
measured geometric law with the same rate for the windowed estimator's
error, over the stated range of $r$. A two-sided bound establishing the
second from the first is left open.

\begin{figure}[t]\centering
\includegraphics[width=\linewidth]{fig_local_vs_full.pdf}
\caption[Geometric decay of the windowed estimator's error]{Locality of the joint score ($\alpha = 0.8$, $L = 121$,
deterministic RMS over the data law). Left: exact truncation error of the
radius-$r$ local estimator (dots) at three diffusion times; the dashed
lines have the predicted slopes $q(\alpha,t)^r$ from
\eqref{eq:g-bulk-cov}, with no fitting. Right: the same error at fixed
radius as a function of $t$: it vanishes at both ends and peaks at
intermediate $t$, where the band of $\Qt$ is fullest.}
\label{fig:g-locality}
\end{figure}

\subsection{Truncating the matrix versus truncating the inference}

A truncated score can be built in two inequivalent ways.
\emph{Truncate the answer}: band the matrix,
$\widehat S = -Q_t^{(b)} x$ with entries beyond $|i-j| \leq b$ zeroed.
\emph{Truncate the algorithm}: let messages travel at most $r$ hops, i.e.\
run exact inference on the window $x_{k-r}, \dots, x_{k+r}$ (the local
estimator above). Both are linear in $x$, so their RMS errors are again
deterministic. Table~\ref{tab:g-truncation} compares them at equal locality
budget ($b = r = 1$, $L = 40$, $\alpha = 0.8$):
\begin{center}
\begin{minipage}{\linewidth}\centering
\captionof{table}[Truncating the answer against truncating the algorithm]{Truncating the answer against truncating the algorithm, at equal locality budget ($b = r = 1$, $L = 40$, $\alpha = 0.8$). Both estimators are linear in $x$, so both errors are deterministic rather than sampled. The cost of locality peaks at intermediate diffusion times and the ordering reverses nowhere.}
\label{tab:g-truncation}
\begin{tabular}{@{}lcccc@{}}
\toprule
relative RMS score error & $t = 0.05$ & $t = 0.3$ & $t = 1$ & $t = 3$ \\
\midrule
banded matrix ($b = 1$)      & $0.210$ & $0.302$ & $0.135$ & $0.0035$ \\
truncated inference ($r = 1$) & $0.123$ & $0.235$ & $0.130$ & $0.0035$ \\
\bottomrule
\end{tabular}
\end{minipage}
\end{center}
Two observations. The cost of locality is non-monotone in $t$ and peaks at
intermediate times, consistent with Section~\ref{sec:g-tridiag-loss}. And
at small $t$ truncating the algorithm beats truncating the answer
($12.3\%$ against $21.0\%$ at $t = 0.05$): the windowed inference solves a
correctly normalised problem on its window, while zeroing entries of an
inverse produces a matrix that is the inverse of nothing and misstates the
diagonal. The two coincide at large $t$, where the off-diagonal content
vanishes. Either error is measured to decay geometrically in the allowed
range at the rate $q(\alpha, t)$ of Proposition~\ref{thm:g-locality}.
